% GENERATED FILE. Edit canonical lessons, then run npm run generate:books.
% canonical-source-hash: fnv1a64:394819f5ff9cacfb
% canonical-lessons: SA-C406-adhirohati, SA-C406-avatarati, SA-C406-sammardah, SA-C406-panktih, SA-C406-avakasah

\chapter{Getting On and Off}
\label{ch:sa-getting-on-and-off}
\hlchaptermodality{406}

\begin{chapteropening}
\textbf{By the end of this chapter:} \emph{I can say how I got on, found room and got off.}

Complete the last lesson of chapter 406: I can say how I got on, found room and got off.
\end{chapteropening}

\section[\texorpdfstring{\sk{अधिरोहति} (adhirohati)}{अधिरोहति (adhirohati)}]{\sk{अधिरोहति} (adhirohati) — gets on; mounts}

\label{lesson:SA-C406-adhirohati}



\begin{quote}
Before the new one: say the Sanskrit for last, then the Sanskrit for arrives.
\end{quote}

\subsection*{You'll want to know: \sk{अधिरोहति}}
\textbf{\sk{अधिरोहति}} — \emph{adhirohati} — \textquotedblleft{}gets on; mounts\textquotedblright{}.

\begin{quote}
\sk{यात्री} \sk{लोकयानम्} \sk{अधिरोहति।}
\end{quote}

\begin{quote}
\emph{yātrī lokayānam adhirohati.}
\end{quote}

\begin{quote}
The passenger gets on the bus.
\end{quote}

Past (he …): \sk{सः} \sk{लोकयानम्} \sk{अधिरूढवान्}, \emph{saḥ lokayānam adhirūḍhavān}.

\begin{cousinweb}
\textbf{\sk{अधिरोहति}} (\emph{adhirohati}) is built from \emph{adhi}, \textquotedblleft{}over, onto\textquotedblright{}, and the root √ruh, \textquotedblleft{}to rise\textquotedblright{}, which also gives \emph{rohati}, \textquotedblleft{}grows\textquotedblright{}; MW gives \emph{adhi-ruh} as to rise above, ascend or mount, and getting on a bus or a train is that sense in modern use (MW).
\end{cousinweb}

\begin{grammarlens}[title={What you've built}]
One more word for getting around.
\end{grammarlens}

\subsection*{Your turn}
\begin{itemize}
\raggedright
  \item \emph{Say it:} \emph{adhirohati}
  \item \emph{Say it:} \emph{adhirohati}, once more
  \item \emph{Say it:} \emph{yātrī lokayānam adhirohati.} — The passenger gets on the bus.
  \item \emph{Recall:} say the Sanskrit for last, then the Sanskrit for arrives, then say \emph{adhirohati} again
\end{itemize}

\begin{tcolorbox}[breakable,colback=teal!4,colframe=teal!35!black,title={Before you move on}]
What does \sk{अधिरोहति} mean? (\textquotedblleft{}Gets on; mounts\textquotedblright{}.) Say it once more.
\end{tcolorbox}
\section[\texorpdfstring{\sk{अवतरति} (avatarati)}{अवतरति (avatarati)}]{\sk{अवतरति} (avatarati) — gets off; alights}

\label{lesson:SA-C406-avatarati}



\begin{quote}
Before the new one: say the Sanskrit for arrives, then the Sanskrit for gets on.
\end{quote}

\subsection*{You'll want to know: \sk{अवतरति}}
\textbf{\sk{अवतरति}} — \emph{avatarati} — \textquotedblleft{}gets off; alights\textquotedblright{}.

\begin{quote}
\sk{यात्री} \sk{स्थानके} \sk{अवतरति।}
\end{quote}

\begin{quote}
\emph{yātrī sthānake avatarati.}
\end{quote}

\begin{quote}
The passenger gets off at the station.
\end{quote}

Past (he …): \sk{सः} \sk{स्थानके} \sk{अवतीर्णवान्}, \emph{saḥ sthānake avatīrṇavān}.

\begin{cousinweb}
\textbf{\sk{अवतरति}} (\emph{avatarati}) is built from \emph{ava}, \textquotedblleft{}down\textquotedblright{}, and the root √tṝ, \textquotedblleft{}to cross\textquotedblright{}; MW gives \emph{ava-tṝ} as to descend or alight, with the participle \emph{avatīrṇa}, \textquotedblleft{}alighted\textquotedblright{} (MW).
\end{cousinweb}

\begin{grammarlens}[title={What you've built}]
One more word for getting around.
\end{grammarlens}

\subsection*{Your turn}
\begin{itemize}
\raggedright
  \item \emph{Say it:} \emph{avatarati}
  \item \emph{Say it:} \emph{avatarati}, once more
  \item \emph{Say it:} \emph{yātrī sthānake avatarati.} — The passenger gets off at the station.
  \item \emph{Recall:} say the Sanskrit for arrives, then the Sanskrit for gets on, then say \emph{avatarati} again
\end{itemize}

\begin{tcolorbox}[breakable,colback=teal!4,colframe=teal!35!black,title={Before you move on}]
What does \sk{अवतरति} mean? (\textquotedblleft{}Gets off; alights\textquotedblright{}.) Say it once more.
\end{tcolorbox}
\section[\texorpdfstring{\sk{सम्मर्दः} (sammardaḥ)}{सम्मर्दः (sammardaḥ)}]{\sk{सम्मर्दः} (sammardaḥ) — a crush of people}

\label{lesson:SA-C406-sammardah}



\begin{quote}
Before the new one: say the Sanskrit for gets on, then the Sanskrit for gets off.
\end{quote}

\subsection*{You'll want to know: \sk{सम्मर्दः}}
\textbf{\sk{सम्मर्दः}} — \emph{sammardaḥ} — \textquotedblleft{}a crush of people\textquotedblright{}.

\begin{quote}
\sk{प्रातः} \sk{लोकयाने} \sk{महान्} \sk{सम्मर्दः} \sk{भवति।}
\end{quote}

\begin{quote}
\emph{prātaḥ lokayāne mahān sammardaḥ bhavati.}
\end{quote}

\begin{quote}
In the morning there is a great crush in the bus.
\end{quote}

A masculine noun.

\begin{cousinweb}
\textbf{\sk{सम्मर्दः}} (\emph{sammardaḥ}) is built from \emph{sam}, \textquotedblleft{}together\textquotedblright{}, and the root √mṛd, \textquotedblleft{}to press, to crush\textquotedblright{}; MW gives \emph{sam-marda} as pressing together and trampling, and \emph{jana-sammarda} as a throng of people; the crush in a bus is that sense in modern use (MW).
\end{cousinweb}

\begin{grammarlens}[title={What you've built}]
One more word for getting around.
\end{grammarlens}

\subsection*{Your turn}
\begin{itemize}
\raggedright
  \item \emph{Say it:} \emph{sammardaḥ}
  \item \emph{Say it:} \emph{sammardaḥ}, once more
  \item \emph{Say it:} \emph{prātaḥ lokayāne mahān sammardaḥ bhavati.} — In the morning there is a great crush in the bus.
  \item \emph{Recall:} say the Sanskrit for gets on, then the Sanskrit for gets off, then say \emph{sammardaḥ} again
\end{itemize}

\begin{tcolorbox}[breakable,colback=teal!4,colframe=teal!35!black,title={Before you move on}]
What does \sk{सम्मर्दः} mean? (\textquotedblleft{}A crush of people\textquotedblright{}.) Say it once more.
\end{tcolorbox}
\section[\texorpdfstring{\sk{पंक्तिः} (paṅktiḥ)}{पंक्तिः (paṅktiḥ)}]{\sk{पंक्तिः} (paṅktiḥ) — a queue; a line of people}

\label{lesson:SA-C406-panktih}



\begin{quote}
Before the new one: say the Sanskrit for gets off, then the Sanskrit for a crush of people.
\end{quote}

\subsection*{You'll want to know: \sk{पंक्तिः}}
\textbf{\sk{पंक्तिः}} — \emph{paṅktiḥ} — \textquotedblleft{}a queue; a line of people\textquotedblright{}.

\begin{quote}
\sk{जनाः} \sk{यानपत्राय} \sk{पंक्तौ} \sk{तिष्ठन्ति।}
\end{quote}

\begin{quote}
\emph{janāḥ yānapatrāya paṅktau tiṣṭhanti.}
\end{quote}

\begin{quote}
People stand in a queue for tickets.
\end{quote}

A feminine noun.

\begin{cousinweb}
\textbf{\sk{पंक्तिः}} (\emph{paṅktiḥ}) is built on \emph{pañcan}, \textquotedblleft{}five\textquotedblright{}: MW gives \emph{paṅkti} first as a set of five, then as a row, a line, a series; a queue is the row in modern use (MW).
\end{cousinweb}

\begin{grammarlens}[title={What you've built}]
One more word for getting around.
\end{grammarlens}

\subsection*{Your turn}
\begin{itemize}
\raggedright
  \item \emph{Say it:} \emph{paṅktiḥ}
  \item \emph{Say it:} \emph{paṅktiḥ}, once more
  \item \emph{Say it:} \emph{janāḥ yānapatrāya paṅktau tiṣṭhanti.} — People stand in a queue for tickets.
  \item \emph{Recall:} say the Sanskrit for gets off, then the Sanskrit for a crush of people, then say \emph{paṅktiḥ} again
\end{itemize}

\begin{tcolorbox}[breakable,colback=teal!4,colframe=teal!35!black,title={Before you move on}]
What does \sk{पंक्तिः} mean? (\textquotedblleft{}A queue; a line of people\textquotedblright{}.) Say it once more.
\end{tcolorbox}
\section[\texorpdfstring{\sk{अवकाशः} (avakāśaḥ)}{अवकाशः (avakāśaḥ)}]{\sk{अवकाशः} (avakāśaḥ) — room; free space}

\label{lesson:SA-C406-avakasah}



\begin{quote}
Before the new one: say the Sanskrit for a crush of people, then the Sanskrit for a queue.
\end{quote}

\subsection*{You'll want to know: \sk{अवकाशः}}
\textbf{\sk{अवकाशः}} — \emph{avakāśaḥ} — \textquotedblleft{}room; free space\textquotedblright{}.

\begin{quote}
\sk{लोकयाने} \sk{अवकाशः} \sk{न} \sk{अस्ति।}
\end{quote}

\begin{quote}
\emph{lokayāne avakāśaḥ na asti.}
\end{quote}

\begin{quote}
There is no room in the bus.
\end{quote}

A masculine noun.

\begin{cousinweb}
\textbf{\sk{अवकाशः}} (\emph{avakāśaḥ}) is built from \emph{ava} and the root √kāś, \textquotedblleft{}to appear, to shine\textquotedblright{}; MW gives \emph{avakāśa} as place, space, room, and also an occasion or opportunity (MW).
\end{cousinweb}

\begin{grammarlens}[title={What you've built}]
One more word for getting around.
\end{grammarlens}

\subsection*{Your turn}
\begin{itemize}
\raggedright
  \item \emph{Say it:} \emph{avakāśaḥ}
  \item \emph{Say it:} \emph{avakāśaḥ}, once more
  \item \emph{Say it:} \emph{lokayāne avakāśaḥ na asti.} — There is no room in the bus.
  \item \emph{Recall:} say the Sanskrit for a crush of people, then the Sanskrit for a queue, then say \emph{avakāśaḥ} again
\end{itemize}

\begin{tcolorbox}[breakable,colback=teal!4,colframe=teal!35!black,title={Before you move on}]
What does \sk{अवकाशः} mean? (\textquotedblleft{}Room; free space\textquotedblright{}.) Say it once more.
\end{tcolorbox}
